\chapter{Conclusion}
\label{chap:conclusion}

This thesis asked what it takes to scale a scene-level 3D Gaussian Splatting variational autoencoder from a few hundred training scenes (the regime of Can3Tok \citep{gao2025can3tok}) to several thousand pre-fitted heterogeneous indoor scenes (the regime made available by SceneSplat-7K \citep{li2025scenesplat}), and what the design implications are for a downstream flow-matching generator. I conclude by returning to the three research questions and listing the next concrete steps.

%% ──────────────────────────────────────────────────────────────────────────
\section{Answers to the Research Questions}
\label{sec:conc_answers}

\paragraph{RQ1.} Yes, imposing a canonical spatial order on the Gaussian set restores convergence at scale, but only when paired with a decoder that can use it. A fixed-frame Hilbert serialisation (\S\ref{sec:method_hilbert}) drops validation $\ell_2$ from 99 to 77 with no other change (V1 $\to$ V2). Adding anchor-relative position decoding on top of the ordering drops it further to 45 (V2 $\to$ V3). The ordering and the anchor decomposition together turn an unlearnable slot-aligned objective into a learnable one.

\paragraph{RQ2.} Yes, a windowed local encoder plus a shared per-token decoder closes most of the train-to-test gap. The two components are interlocking: the local encoder needs the token-local decoder to consume per-region features cleanly, and the token-local decoder needs the local encoder to feed it the right kind of context. Together, the decoder parameter count drops from about 777M to about 1.6M (a $500\times$ reduction) while validation reconstruction reaches a level comparable to the global-decoder run.

\paragraph{RQ3.} Yes, per-Gaussian semantic supervision and a layout/geometry latent split can be added without degrading reconstruction. The 300-scene exploratory study (Strategy~B1 with Pool+pgNCE-style InfoNCE, \S\ref{sec:results_300}) confirms that per-Gaussian features cluster by ScanNet-72 category and that a deterministic layout code is stable under partial-scan masking. At scale, the V6/V7 variants match V3 on reconstruction while producing the split-schema latent that the hierarchical Stage~2 generator (\S\ref{sec:method_stage2}) needs.

\paragraph{Main question.} The combination of fixed-frame Hilbert serialisation, anchor-relative position decoding, and a local-encoder plus token-local-decoder pair is what scales the Stage~1 VAE from hundreds to several thousand pre-fitted indoor scenes. Each ingredient on its own is insufficient; the three together turn an architecture that fails to converge at scale into one that reconstructs held-out scenes at a useful level with a 500$\times$ smaller decoder. Layered on top, a layout/geometry split with per-Gaussian semantic supervision is a viable schema for the downstream generator.

%% ──────────────────────────────────────────────────────────────────────────
\section{Future Work}
\label{sec:conc_future}

The next milestones for the project follow directly from the limitations in \S\ref{sec:disc_limitations}.

\begin{enumerate}
\item \textbf{Finish Stage~2 end-to-end.} Generate samples from the trained Stage~2 models for variants 5, 6, and 7, render them, and report PSNR, SSIM, and LPIPS against held-out ground truth. The same checkpoints support the completion task: mask 30--80\% of the input Hilbert blocks of a held-out scene and recover the masked regions through Stage~2.
\item \textbf{Density-uniform input sampling.} Replace the top-opacity ranking of the 40{,}000 input Gaussians with farthest-point sampling on positions, weighted by opacity. This should give a more spatially uniform input distribution and reduce the bias towards opaque furniture surfaces.
\item \textbf{Rendering-quality evaluation throughout.} Add rendered-image metrics (PSNR/SSIM/LPIPS) on the reconstructed scenes for all seven variants, not just the generative outputs. This gives the perceptual companion metric to the per-primitive $\ell_2$.
\item \textbf{Unified label space.} Bring ARKitScenes and ScanNet++ under a unified per-Gaussian taxonomy (a multi-taxonomy classification head, with per-dataset masks) so the InfoNCE supervision can be active on the full training corpus.
\item \textbf{Outdoor scenes.} Test the same architecture on outdoor 3DGS data with a contracted-depth normalisation. The Hilbert serialisation and the anchor scaffold should both generalise, but the canonical extent has to change.
\item \textbf{Joint-splat ablation.} Re-train one of the Stage~1 variants jointly with the Gaussian fitting at scale and compare to the same variant on the pre-fitted SceneSplat-7K splats. This is the cleanest test of the frozen-splat hypothesis (\S\ref{sec:disc_frozen}).
\item \textbf{Multiple seeds.} Re-run V3, V5, V6, V7 on at least three seeds each to get variance estimates on the validation $\ell_2$ comparisons.
\end{enumerate}

The codebase, the SLURM job scripts, the YAML configs for each of the seven variants, and the Stage~2 model code are available at \repo. The Stage~1 checkpoints used to produce Table~\ref{tab:7_summary} are released alongside, so the next person can pick up directly from the Stage~2 milestone without retraining.